\documentclass[11pt]{article}

\usepackage[margin=1in]{geometry}
\usepackage{booktabs}
\usepackage{multirow}
\usepackage{amsmath}
\usepackage{natbib}
\usepackage{xcolor}

% Skeleton conventions:
%  - Every section carries \itemize bullets = the meat to be expanded into prose.
%  - \todo{...} = decision or number still needed before writing the prose.
%  - Single-table constraint: tab:main-results is the ONLY table in the paper.
%    All ablation/analysis numbers go into prose deltas or figures.
\newcommand{\todo}[1]{\textcolor{red}{[TODO: #1]}}
\newcommand{\method}{SCOTECI}

\title{\method{}: Synthetic Chain-of-Thought Demonstrations for Document-Level Event Causality Identification\\
{\normalsize \todo{working title --- alternatives: ``What Matters in Document-Scale Few-Shot ECI?'' / ``Rationalize, Select, Resample: Demonstration Design for Document-Level ECI''}}}
\author{}
\date{}

\begin{document}

\maketitle

% ======================================================================
\begin{abstract}
\begin{itemize}
  \item LLMs applied to Event Causality Identification (ECI) suffer from causal
        hallucination: near-total recall, unusable precision \citep{gao2023chatgpt}.
  \item We propose \method{}: LLM-based ECI at the \emph{document} level --- the
        full document text sits in the context and one call labels every pair
        jointly under a structured reasoning protocol with explicit coherence
        rules; demonstrations are whole labelled documents equipped with
        \emph{synthetic chain-of-thought}: gold-guided reasoning generated by an
        LLM, then decontaminated so it reads as genuine discovery.
  \item Beyond the method, a systematic study of the demonstration design
        space: (i) CoT synthesis protocol (2-step gold-guided vs.\ 3-step
        blind-attempt-then-correct); (ii) demonstration selection (random,
        constant, TF-IDF, sentence-embedding, mention-set overlap);
        (iii) resampling with majority vote --- number of passes and whether
        each pass sees the \emph{same} or a \emph{distinct} demonstration set;
        (iv) backbone choice across four contemporary LLMs.
  \item No fine-tuning, no human rationales, no external knowledge base.
  \item Headline: demonstration design dominates backbone choice --- selection
        strategy and CoT protocol move F1 by \todo{X} points where swapping
        backbones moves it by \todo{Y}; on Causal-TimeBank, F1
        \todo{re-run on final backbone ladder} with balanced precision/recall
        instead of the degenerate high-recall regime.
\end{itemize}
\end{abstract}

% ======================================================================
\section{Introduction}
\label{sec:intro}

\begin{itemize}
  \item \textbf{Task and stakes.} ECI = given a document and event mentions, decide for
        event pairs whether a (directed) causal relation holds
        \citep{cheng2025survey}. Feeds event forecasting, QA, narrative understanding.
  \item \textbf{The failure mode.} Prompted LLMs need no training data but
        exhibit \emph{causal hallucination} \citep{gao2023chatgpt}: recall
        of 80--96\% with precision below 10\% on sparse corpora --- the model
        asserts causality wherever two events co-occur.
  \item \textbf{Idea 1 --- operate at the document scale.} The full document
        text is placed in the context window and a single call labels
        \emph{all} candidate pairs jointly, under an in-prompt protocol
        (mention pre-filter $\to$ per-pair for/against arguments $\to$ plausibility
        ranking $\to$ final JSON) plus explicit coherence rules (symmetry,
        dataset-specific transitivity). The model sees the whole narrative and
        produces the document's causal graph in one coherent generation ---
        deliberately betting on long-context reasoning quality, so the method
        \emph{improves as models improve}.
  \item \textbf{Idea 2 --- synthetic chain-of-thought few-shots.} Document-level
        demonstrations are whole labelled documents from the training split.
        Bare input$\to$gold-JSON demos teach the output format, not
        the reasoning; human-written document-level rationales do not exist.
        We therefore \emph{generate} them:
        \begin{itemize}
          \item gold-guided generation: the demo document is presented under the
                exact task prompt together with its gold labels; the LLM writes
                reasoning that leads naturally to those labels;
          \item decontamination: the reasoning is rewritten from scratch to
                purge any foreknowledge phrasing, so every sentence reads as
                discovery, not confirmation;
          \item 3-step variant: blind attempt $\to$ gold-based
                correction $\to$ decontamination (preserves the texture of a
                genuine attempt);
          \item injected as multi-turn (user, assistant) exchanges; cached per
                demo document.
        \end{itemize}
  \item \textbf{Idea 3 --- the demonstration design space is the lever.} With
        the formulation fixed, three orthogonal knobs remain: \emph{how} the
        synthetic reasoning is produced (2- vs.\ 3-step), \emph{which}
        documents are shown (five selection strategies from query-independent
        to semantically retrieved), and \emph{how many} votes are taken with
        \emph{what} demonstration diversity (resampling with shared vs.\
        distinct demonstration sets per pass). We map this space systematically
        across four backbones.
  \item \textbf{Results preview.} Single evaluation table
        (Table~\ref{tab:main-results}): \todo{headline numbers on the final
        backbone ladder --- pilot runs give CTB F1 51.0 with P/R 53.2/49.0,
        i.e.\ the balanced regime, vs.\ single-digit precision for zero-shot
        prompting}. Ablations: \todo{selection-strategy and CoT-depth deltas};
        backbone choice is \emph{not} the dominant factor \todo{quantify}.
  \item \textbf{Contributions.}
        \begin{itemize}
          \item A document-level formulation of LLM-based ECI --- full text in
                context, one generation labels every pair, structured protocol
                and in-prompt coherence rules, $O(1)$ calls per document.
          \item A synthetic-CoT demonstration pipeline (gold-guided generation +
                decontamination, with a 3-step blind-then-correct variant),
                model-agnostic, requiring no human rationales and no knowledge
                base.
          \item A systematic study of demonstration selection, CoT synthesis
                depth, and demonstration-varied resampling, showing that
                demonstration design outweighs backbone choice
                \todo{confirm once the grid lands}.
        \end{itemize}
\end{itemize}

% ======================================================================
\section{Related Work}
\label{sec:related}

\subsection{LLMs and Causal Reasoning over Text}
\begin{itemize}
  \item \citet{gao2023chatgpt}: systematic evaluation; ChatGPT-class models
        over-predict causality (causal hallucination); established the
        preprocessing and evaluation protocol we follow.
  \item Surveys position ECI as a core event-semantics task with growing
        LLM interest \citep{cheng2025survey}.
  \item Gap: prompting work reports either zero-shot or label-only few-shot
        results; the demonstration design space --- what the demonstrations
        contain, how they are chosen, how they interact with resampling ---
        is unmapped at document scale. That map is this paper.
\end{itemize}

\subsection{In-Context Learning and Demonstration Selection}
\begin{itemize}
  \item ICL emerges with scale \citep{brown2020language} and is highly
        sensitive to demonstration choice and content
        \citep{liu2022makes, min2022rethinking}.
  \item Retrieval-based selection (lexical or embedding similarity to the
        query) is the standard recipe \citep{liu2022makes}; whether it holds
        when each demonstration is an entire labelled document --- long,
        expensive, and few --- is one of our questions
        (\S\ref{sec:ablations-selection}).
\end{itemize}

\subsection{Chain-of-Thought, Synthetic Rationales, and Resampling}
\begin{itemize}
  \item CoT prompting boosts multi-step reasoning \citep{wei2022chain};
        Auto-CoT generates rationales with zero-shot-CoT over clustered queries
        \citep{zhang2023autocot}; STaR bootstraps rationales via
        \emph{rationalization} --- generating reasoning conditioned on the known
        answer --- but consumes them as fine-tuning data \citep{zelikman2022star};
        step-by-step distillation extracts rationales as training signal
        \citep{hsieh2023distilling}.
  \item Self-consistency samples multiple reasoning paths and majority-votes
        the answers \citep{wang2023selfconsistency}; the paths vary by sampling
        temperature. Our resampling instead varies the \emph{evidence}: at
        temperature 0, each pass can receive a different pair ordering and,
        optionally, a different demonstration set --- prompt-side rather than
        decoder-side diversity.
  \item Ours differs: rationalization used \emph{at inference time} as ICL
        demonstrations, at document scale, with an explicit decontamination pass
        to remove answer-leakage phrasing --- without which the demonstrations
        teach the model to assert rather than reason.
\end{itemize}

% ======================================================================
\section{Task Formulation}
\label{sec:task}

\begin{itemize}
  \item Input: document $D$ with inline-marked event mentions
        $M = \{m_1, \dots, m_n\}$ (\texttt{<ID word>} format), and a pair set
        $P$: the annotated candidate list when the corpus provides one
        (constrained setting), otherwise all unordered mention pairs
        (open-ended setting).
  \item Output: a complete labelling $\ell : P \to \{\textit{causes},
        \textit{causedby}, \textit{norel}\}$ --- a unified directed label space
        into which each corpus' native labels are mapped
        \todo{spell out per-dataset label mapping (CAUSE/PRECONDITION,
        FALLING\_ACTION, CauseEffect/EffectCause)}.
  \item Document-scale inference: the \emph{entire} document text is in the
        context --- no windowing, no chunking --- and the model produces $\ell$
        for \emph{all} of $P$ in a single generation, as a JSON array. Every
        labelling decision is made with the full narrative visible, so
        cross-sentence chains, confounders, and discourse structure are
        available to every pair judgement simultaneously.
  \item Evaluation: micro precision/recall/F1 over the positive (causal)
        classes; \textit{norel} excluded; intra- vs.\ inter-sentence subsets
        where the corpus distinguishes them; preprocessing follows
        \citet{gao2023chatgpt}.
\end{itemize}

% ======================================================================
\section{Method}
\label{sec:method}

\subsection{Overview}
\begin{itemize}
  \item Pipeline (Figure~\ref{fig:overview}): test document $\to$ select $k$
        demonstration documents from the training split $\to$ generate (and
        cache) synthetic CoT for each $\to$ assemble multi-turn prompt
        (demonstrations as past exchanges, then the test document) $\to$
        document-level LLM call(s) $\to$ parse JSON $\to$ optional resampling
        with majority vote.
  \item Four components, each carrying a studied design dimension: the
        document-level structured prompt (\S\ref{sec:docprompt}), demonstration
        selection (\S\ref{sec:selection}), synthetic-CoT generation with its
        2- and 3-step protocols (\S\ref{sec:synthcot}), and resampling with
        shared or distinct demonstration sets (\S\ref{sec:inference}).
\end{itemize}

% \begin{figure}[t]
%   \centering
%   \todo{Figure 1: pipeline overview --- demonstration selection, 2-/3-step CoT
%   synthesis (gold-guided generation -> decontamination), multi-turn prompt
%   assembly, doc-level call(s), majority vote over resamples with shared or
%   distinct demonstration sets.}
%   \caption{Overview of \method{}.}
%   \label{fig:overview}
% \end{figure}

\subsection{Document-Level Structured Prompting}
\label{sec:docprompt}
\begin{itemize}
  \item One call per document; the full text in context; every pair of $P$
        requested at once; the model sees the whole narrative, not a windowed
        snippet.
  \item In-prompt reasoning protocol (four ordered steps, enforced by the
        system prompt):
        \begin{itemize}
          \item Step 1 --- mention-level pre-filter: affirmatively argued
                exclusions (``$m$ is a location with no agentive role''), never
                blanket dismissal; surviving mentions listed.
          \item Step 2 --- per-pair arguments \emph{for} (connective, mechanism,
                counterfactual asymmetry, discourse structure) and
                \emph{against} (temporal-only, confounder, reporting verb,
                distance).
          \item Step 3 --- qualitative ranking of retained pairs by causal
                plausibility, no labels committed yet.
          \item Step 4 --- final JSON over \emph{every} requested pair;
                everything pre-filtered defaults to \textit{norel}.
        \end{itemize}
  \item Precision-first rule block (targets causal hallucination directly):
        two-signal rule for non-overt links; stricter gate for inter-sentential
        / long-range pairs; default-to-\textit{norel} under uncertainty;
        confounder / proximal-cause check; reporting and temporal connectives
        explicitly ruled non-causal; counterfactual direction test.
  \item In-prompt coherence rules: symmetry
        ($\textit{causes}(i,j) \Leftrightarrow \textit{causedby}(j,i)$) and
        dataset-specific transitivity rules --- checkable at parse time because
        the whole graph is produced jointly \todo{decide whether the coherence
        checker/post-hoc rule pruning appears in main text or appendix}.
  \item Why this leverages better models: the protocol asks for exactly the
        global, long-context reasoning that frontier models are differentially
        good at.
\end{itemize}

\subsection{Demonstration Selection}
\label{sec:selection}
\begin{itemize}
  \item Demonstration pool: training split of the same corpus (capped pool
        size to bound CoT-generation cost) \todo{final pool size}.
  \item Demonstrations are \emph{whole documents with their complete gold
        labelling}, including the \textit{norel} pairs --- every demo carries
        explicit negative evidence alongside the positives.
  \item Five selection strategies, spanning query-independent to
        semantically adaptive (all evaluated in
        \S\ref{sec:ablations-selection}):
        \begin{itemize}
          \item \emph{random} --- $k$ documents drawn uniformly per test
                document; query-independent and varying;
          \item \emph{constant} --- one fixed set of $k$ documents reused for
                every test document (chosen once, e.g.\ the most compact
                well-annotated documents in the pool); query-independent and
                frozen --- isolates whether \emph{adaptivity} matters at all,
                and maximizes prompt-cache reuse;
          \item \emph{TF-IDF} --- cosine over full document text, top-$k$;
                lexical similarity;
          \item \emph{sentence-embedding} --- cosine over pooled
                sentence-transformer embeddings of the document text, top-$k$;
                semantic similarity \todo{state embedding model};
          \item \emph{mention-set overlap} --- Jaccard similarity between the
                \emph{event mention sets} of test and candidate documents:
                two documents that share many event types (``attack'',
                ``collapse'', ``rescue'') likely exercise similar causal
                schemas, independently of surface wording.
        \end{itemize}
  \item Default configuration: TF-IDF, $k=3$ \todo{confirm after the
        selection ablation}.
\end{itemize}

\subsection{Synthetic Chain-of-Thought Generation}
\label{sec:synthcot}
\begin{itemize}
  \item Motivation: label-only demos teach the output schema, not the
        reasoning; document-level human rationales are unavailable and
        prohibitively expensive to write; gold labels for training documents
        are available --- so rationales can be \emph{synthesized}.
  \item 2-step protocol, per demonstration document:
        \begin{itemize}
          \item (i) \emph{gold-guided generation}: the demo document is shown
                under the exact task prompt, plus its gold labels and the
                instruction to write reasoning that ``leads naturally to these
                labels from the text alone'' without revealing foreknowledge;
          \item (ii) \emph{decontamination}: the response is rewritten from
                scratch to purge leakage phrasing (``the labels show\dots'',
                ``as indicated by the expected output\dots''), so each sentence
                reads as analysis, not confirmation.
        \end{itemize}
  \item 3-step protocol: (i) \emph{blind attempt} with no gold access ---
        the model reasons over the demo document exactly as it would at test
        time; (ii) \emph{gold-based correction} that fixes wrong conclusions
        while keeping correct reasoning; (iii) decontamination.
        Hypothesis: the blind attempt preserves the texture, hedging, and
        difficulty profile of genuine inference --- reasoning the test-time
        model can actually imitate --- where gold-guided generation may read
        as too fluent, arguing every pair with unearned confidence.
        Costs one extra generation call per demo. The 2- vs.\ 3-step
        comparison is a primary ablation (\S\ref{sec:ablations-cot}).
  \item Injection: each demo becomes a (user task turn, assistant CoT turn)
        exchange preceding the test document --- the model ``remembers having
        done'' $k$ documents correctly, reasoning included; compatible with
        multi-step task prompts (one synthetic response per step).
  \item Generated once per (document, protocol), cached across the run and
        on disk across runs; generator model = inference model by default
        (self-generated rationales) \todo{teacher-vs-self generation
        experiment --- does a stronger generator lift a weaker reasoner?}.
  \item Relation to STaR-style rationalization \citep{zelikman2022star}:
        same trick --- condition on the answer to obtain reasoning --- but used
        as in-context demonstrations rather than fine-tuning data, and with an
        explicit decontamination pass, which we show is load-bearing
        (\S\ref{sec:ablations-cot}).
\end{itemize}

\subsection{Resampling and Aggregation}
\label{sec:inference}
\begin{itemize}
  \item Temperature 0; strict JSON parsing with bounded retries; missing pairs
        default to \textit{norel}.
  \item Resampling: $N$ independent passes with per-pair majority vote, ties
        broken to \textit{norel} (precision-preserving). Because decoding is
        greedy, diversity is injected on the \emph{prompt side}:
        \begin{itemize}
          \item \emph{pair-order reshuffle}: the candidate pair list is
                re-ordered per pass (deterministically seeded per document and
                pass index, so runs are reproducible) --- breaks position bias;
          \item \emph{shared demonstrations} (baseline): every pass sees the
                same $k$ demos --- votes differ only through pair order;
          \item \emph{distinct demonstrations}: $N \times k$ documents are
                selected by the active strategy and split round-robin across
                passes (each pass gets a rank-balanced set of $k$) --- every
                vote is grounded in different evidence, decorrelating errors
                that stem from an unlucky demonstration draw. Requires a pool
                of at least $N \times k$ eligible documents.
        \end{itemize}
  \item The resampling study (\S\ref{sec:ablations-resampling}) crosses
        $N \in \{1, 3, 5\}$ with \{shared, distinct\} demonstrations:
        does voting help, and does demonstration diversity add anything beyond
        pair-order diversity? \todo{run the grid}.
  \item Cost: $N$ calls per document plus $2$--$3k$ amortized CoT-generation
        calls (cached); distinct demonstrations multiply CoT generation by
        $N$ but not test-time calls; quantified in \S\ref{sec:cost}.
\end{itemize}

% ======================================================================
\section{Experiments}
\label{sec:experiments}

\subsection{Datasets and Metrics}
\begin{itemize}
  \item EventStoryLine (ESC) \citep{caselli2017event}: 22 topics, 258 docs,
        1{,}770 causal pairs \todo{confirm version used (0.9 vs 1.5) and align
        statement with Gao preprocessing}.
  \item Causal-TimeBank (CTB) \citep{mirza2014annotating}: 183 docs, 318 causal
        pairs --- sparsest positives, hence the dataset where hallucination is
        most punishing; no intra/inter split.
  \item MAVEN-ERE \citep{wang2022mavenere}: large-scale Wikipedia corpus,
        57{,}992 causal pairs \todo{state evaluated subset/sample size}.
  \item Preprocessing follows \citet{gao2023chatgpt}; micro P/R/F1 on the causal
        classes, \textit{norel} excluded.
  \item \todo{decide whether the MECI multilingual evaluation
        \citep{lai2022meci} goes in (appendix?) --- pipeline supports it.}
\end{itemize}

\subsection{Backbones and Baselines}
\begin{itemize}
  \item Four contemporary backbones spanning tiers and architectures, all via
        a unified API \todo{exact endpoint identifiers}:
        \begin{itemize}
          \item GPT-5.5 (frontier closed);
          \item GLM-5.2 (large open);
          \item DeepSeek-V4 (large open MoE);
          \item Qwen3-30B-Thinking (mid-size open, explicit reasoning mode).
        \end{itemize}
  \item Baselines, same backbone and same document-scale setting so every
        comparison is within-formulation:
        \begin{itemize}
          \item \emph{Base}: task description + document + pair list, direct
                JSON answer;
          \item \emph{CoT} \citep{wei2022chain}: step-by-step reasoning
                requested before the JSON;
          \item \method{}: full pipeline (structured protocol + synthetic-CoT
                demonstrations + resampling) \todo{state final config used in
                the main table}.
        \end{itemize}
  \item Model comparison is a \emph{secondary} axis by design: the paper's
        claim is that demonstration design, not backbone choice, is the
        dominant lever \todo{verify the claim survives the grid}.
\end{itemize}

\subsection{Implementation Details}
\begin{itemize}
  \item $k=3$ demonstrations, TF-IDF selection, 2-step CoT protocol,
        temperature 0, resampling $N=3$ with distinct demonstration sets
        \todo{confirm final default config; pool size; exact prompt files
        versioned in the repo}.
  \item Synthetic CoT generation cost: 2 (or 3) extra calls per unique
        demo document, amortized by caching \todo{give measured token counts}.
  \item Deterministic pipeline: fixed evaluation subset (first $n$ documents
        of each split), per-document seeded pair-order shuffling, per-(document,
        pass) seeded resample reshuffling --- identical runs are exactly
        reproducible up to provider nondeterminism.
\end{itemize}

\subsection{Main Results}
\label{sec:results}

% ----------------------------------------------------------------------
% Single merged results table: 4 backbones x {Base, CoT, SCOTECI} on the
% three datasets. All rows run by us, same document-scale setting.
% Pilot numbers from the DeepSeek-V3.2 development runs, for reference
% while the final grid is pending (do not ship):
%   SCOTECI  ESC 34.7/57.2/43.2 (intra only)  CTB 53.2/49.0/51.0
%            MAVEN 35.5/37.7/36.6 (intra only)
% Bold = best score per column once numbers land.
% ----------------------------------------------------------------------
\begin{table}[htbp]
\centering
\small
\caption{Main results: micro P/R/F1 on the causal classes. All rows use the
same document-scale setting (full text in context, all pairs labelled in one
generation); \method{} adds synthetic-CoT demonstrations and resampling.
\todo{fill after the final grid; mark best per column in bold; state
resampling config; note any intra-sentence-only caveats}.}
\label{tab:main-results}
\begin{tabular}{l ccc|ccc|ccc}
\toprule
 & \multicolumn{3}{c|}{ESC} & \multicolumn{3}{c|}{CTB} & \multicolumn{3}{c}{MAVEN-ERE} \\
Method & P & R & F1 & P & R & F1 & P & R & F1 \\
\midrule
\textbf{GPT-5.5:} & & & & & & & & & \\
\quad Base          & -- & -- & -- & -- & -- & -- & -- & -- & -- \\
\quad CoT           & -- & -- & -- & -- & -- & -- & -- & -- & -- \\
\quad \method{}     & -- & -- & -- & -- & -- & -- & -- & -- & -- \\
\midrule
\textbf{GLM-5.2:} & & & & & & & & & \\
\quad Base          & -- & -- & -- & -- & -- & -- & -- & -- & -- \\
\quad CoT           & -- & -- & -- & -- & -- & -- & -- & -- & -- \\
\quad \method{}     & -- & -- & -- & -- & -- & -- & -- & -- & -- \\
\midrule
\textbf{DeepSeek-V4:} & & & & & & & & & \\
\quad Base          & -- & -- & -- & -- & -- & -- & -- & -- & -- \\
\quad CoT           & -- & -- & -- & -- & -- & -- & -- & -- & -- \\
\quad \method{}     & -- & -- & -- & -- & -- & -- & -- & -- & -- \\
\midrule
\textbf{Qwen3-30B-Thinking:} & & & & & & & & & \\
\quad Base          & -- & -- & -- & -- & -- & -- & -- & -- & -- \\
\quad CoT           & -- & -- & -- & -- & -- & -- & -- & -- & -- \\
\quad \method{}     & -- & -- & -- & -- & -- & -- & -- & -- & -- \\
\bottomrule
\end{tabular}
\end{table}

\begin{itemize}
  \item \textbf{CTB (expected headline).} Pilot runs put \method{} in the
        balanced regime (P $\approx$ R $\approx$ 50) where zero-shot prompting
        is degenerate (single-digit precision at 80--96 recall,
        \citealp{gao2023chatgpt}) \todo{final numbers per backbone}.
  \item \textbf{Reading across method rows.} Base/CoT expected in the
        high-recall/low-precision regime; \method{} trades recall for large
        precision gains --- quantify the P/R rebalancing per dataset
        \todo{once the grid lands}.
  \item \textbf{Reading across backbone blocks.} The between-backbone spread
        within a method row vs.\ the between-method spread within a backbone
        block --- the central comparison for the ``demonstration design
        dominates backbone choice'' claim \todo{compute both spreads}.
  \item \todo{add significance / resampling variance statement once final runs
        land.}
\end{itemize}

% ======================================================================
\section{Analysis}
\label{sec:analysis}
% Single-table constraint: everything below reported as prose deltas
% ("removing X costs Y F1 on CTB") or as figures, never new tables.

\subsection{Synthetic CoT: Protocol Depth and Decontamination}
\label{sec:ablations-cot}
\begin{itemize}
  \item Remove synthetic CoT entirely (demos = document + gold JSON only):
        isolates the contribution of \emph{reasoning-bearing} demonstrations
        \todo{expected: format retained, precision drops --- run}.
  \item 2-step (gold-guided $\to$ decontaminate) vs.\ 3-step (blind attempt
        $\to$ correct $\to$ decontaminate): does preserving the texture of a
        genuine attempt beat fluent gold-guided reasoning? Does the answer
        depend on the backbone (a stronger model may produce blind attempts
        worth correcting; a weaker one may not)? \todo{run both protocols
        across the backbone ladder --- this is the headline ablation}.
  \item Decontamination on/off: does leakage phrasing in demos teach the model
        to assert without reasoning? \todo{run --- this justifies the step}.
  \item Qualitative: side-by-side of a 2-step and a 3-step rationale for the
        same demo document \todo{pick a CTB doc; appendix if long}.
\end{itemize}

\subsection{Demonstration Selection}
\label{sec:ablations-selection}
\begin{itemize}
  \item Grid: \{random, constant, TF-IDF, sentence-embedding, mention-set
        overlap\} $\times$ datasets, default config otherwise \todo{run}.
  \item Hypotheses to test:
        \begin{itemize}
          \item adaptivity gap: do query-dependent strategies (TF-IDF,
                embedding, mention overlap) beat query-independent ones
                (random, constant) when each demo is an entire document?
          \item lexical vs.\ semantic vs.\ event-level: mention-set overlap
                matches on \emph{event vocabulary} rather than surface text ---
                if causal schemas transfer along event types, it should win on
                corpora with recurring scenarios;
          \item constant as the cache-friendly floor: if constant is within
                noise of retrieval, the entire selection machinery can be
                dropped for a fixed cached prefix \todo{cost implications in
                \S\ref{sec:cost}}.
        \end{itemize}
  \item Number of demonstrations $k \in \{1, 3, 5\}$ \todo{run; doc-level
        demos are long, context dilution may bite sooner than in
        short-form ICL}.
\end{itemize}

\subsection{Resampling: Vote Count and Demonstration Diversity}
\label{sec:ablations-resampling}
\begin{itemize}
  \item Grid: $N \in \{1, 3, 5\}$ $\times$ \{shared, distinct\} demonstration
        sets, default config otherwise \todo{run}.
  \item Questions:
        \begin{itemize}
          \item marginal value of votes: F1 vs.\ $N$ --- where does majority
                vote saturate at temperature 0, given that pair-order
                reshuffling is the only variance source in the shared setting?
          \item does demonstration diversity add signal? Distinct sets
                decorrelate demonstration-induced errors: an unlucky demo hurts
                one vote, not all $N$. Expected to matter most for
                high-variance selection (random) and least for constant
                \todo{test the interaction};
          \item precision mechanics: ties break to \textit{norel}, so voting
                is precision-preserving by construction --- decompose the F1
                movement into P and R components \todo{compute}.
        \end{itemize}
  \item Practical note: distinct demonstrations need $N \times k$ eligible
        pool documents and multiply CoT-generation (not test-time) cost by up
        to $N$ \todo{measured overhead}.
\end{itemize}

\subsection{Backbone Comparison}
\label{sec:ablations-models}
\begin{itemize}
  \item Deliberately last: rank the four backbones under the full method and
        under Base --- if the method compresses the between-model gap, the
        demonstrations are doing work the raw model quality otherwise would
        \todo{compute; figure F1-per-backbone under Base vs.\ \method{},
        not a table}.
  \item Reasoning-mode interaction: Qwen3-30B-Thinking emits its own long
        reasoning --- do synthetic-CoT demonstrations compound with or
        substitute for native thinking? \todo{inspect traces; brief prose}.
  \item Cross-model demo reuse: demos generated by one backbone consumed by
        another \todo{optional; run only if time permits}.
\end{itemize}

\subsection{Cost Analysis}
\label{sec:cost}
\begin{itemize}
  \item LLM calls per document: $N$ test-time calls plus $2$--$3$
        CoT-generation calls per unique demo document (cached across the run
        and across passes in the shared setting) \todo{measured token counts;
        cost-vs-F1 figure across configs --- constant selection + shared demos
        as the cheap end, embedding selection + distinct demos as the
        expensive end}.
\end{itemize}

\subsection{Error Analysis}
\begin{itemize}
  \item False-positive taxonomy: transitive flattening (A$\to$B$\to$C read as
        A$\to$C when the scheme does not credit it), temporal-as-causal,
        reporting verbs \todo{sample and count from traces}.
  \item Dense documents: does precision degrade with the number of gold pairs
        per document? \todo{scatter/figure}.
  \item Direction errors: \textit{causes} vs.\ \textit{causedby} confusions
        as a fraction of true-pair errors \todo{compute from logs}.
  \item Qualitative: one worked example --- synthetic CoT demo, model reasoning,
        prediction vs.\ gold \todo{pick a CTB doc}.
\end{itemize}

% ======================================================================
\section{Conclusion}
\begin{itemize}
  \item Formulated LLM-based ECI at the document scale: full text in context,
        a single structured generation labels every pair, coherence rules
        enforced in-prompt.
  \item Introduced synthetic, decontaminated chain-of-thought demonstrations:
        rationalization at inference time, no human rationales, no knowledge
        base, no fine-tuning --- with a 2-step and a 3-step synthesis protocol.
  \item Mapped the demonstration design space --- synthesis depth, five
        selection strategies, resampling count $\times$ demonstration
        diversity --- across four backbones; demonstration design is the
        dominant lever, backbone choice secondary \todo{confirm}.
  \item Large precision-driven gains where hallucination is costliest
        \todo{final CTB numbers}; balanced P/R elsewhere.
  \item Future work: multilingual MECI, per-document strategy routing,
        teacher-generated CoT for small models.
\end{itemize}

% ======================================================================
\section*{Limitations}
\begin{itemize}
  \item Some results are currently intra-sentence only (ESC, MAVEN-ERE);
        full-document numbers must land before the cross-dataset story is
        complete \todo{track}.
  \item Synthetic CoT costs $2$--$3$ generation calls per demonstration
        document (amortized by caching, but nonzero); distinct-demonstration
        resampling multiplies the number of unique demos by up to $N$.
  \item Decontamination is itself LLM-based; residual leakage phrasing cannot
        be ruled out \todo{manual audit of $n$ sampled demos}.
  \item Long documents with large pair sets stress the single-call format
        (output-length limits, attention dilution); per-document routing or
        chunking may be needed at the extreme.
  \item Provider-side nondeterminism means temperature-0 runs are reproducible
        in inputs but not bit-identical in outputs.
\end{itemize}

% ======================================================================
% \appendix
% \section{Prompts}
% \begin{itemize}
%   \item Full system/user prompts per dataset (versioned in repo).
%   \item CoT-generation prompts: gold hint, blind-attempt correction hint,
%         decontamination.
%   \item Coherence rule sets per dataset.
% \end{itemize}
% \section{Additional Details}
% \begin{itemize}
%   \item Dataset statistics and label mappings.
%   \item Resampling/tie-breaking details; JSON-repair procedure.
%   \item Demonstration-set splitting (round-robin over selection ranking).
%   \item MECI multilingual results (if not in main text).
% \end{itemize}

\bibliographystyle{plainnat}
\bibliography{references}

\end{document}
